% Phase 3 CP3c. Control line: does predictive selection raise satisfaction? Table 2 = tables/tab_control.tex.
\subsection{Constraint Satisfaction under Predictive Control}
\label{sec:exp_control}

For RQ2a and RQ2b, we ask whether predictive control using the learned operator improves constraint satisfaction over the same frozen model without control, and how it compares with two white-box latent controllers. The setup, benchmarks, and baselines are described in \S\ref{sec:exp_setup}; Table~\ref{tab:control} reports the constraint-satisfaction score defined there.

Averaged across the eleven subtasks, predictive control improves constraint satisfaction over the uncontrolled baseline by 14.5 points on Qwen3-4B and 14.8 points on Qwen3-8B. These are point estimates, with no seeds or intervals. The controller scores each candidate by its predicted trajectory rather than its immediate effect, allowing it to choose a command different from the target while the attribute is still changing.
Predictive control also outperforms the stronger white-box baseline, RE-Control, by 5.3 points on Qwen3-4B and 5.7 points on Qwen3-8B. These are also point estimates without intervals. RE-Control uses model activations, whereas predictive control operates from text alone.

Across the 22 model-by-subtask comparisons, predictive control scores higher than the uncontrolled baseline in every case and higher than both white-box controllers in every case. This count answers RQ2a and RQ2b affirmatively; it is not a claim of statistical significance.

